%% 第 3 章：特征值、特征向量与特征空间
%% 本文件由 tools/md2tex.py 自动生成，请勿直接编辑。

\chapter{特征值、特征向量与特征空间}


\begin{jdefn}{特征值与特征向量}{r:4:1}
若存在非零向量 $v\in V$，使：

\[
T(v)=\lambda v,
\]

则称 $\lambda\in\mathbb F$ 是 $T$ 的特征值，$v$ 是对应的特征向量。

特征向量张成的一维空间：

\[
\operatorname{span}(v)
\]

是最简单的不变子空间。
\end{jdefn}

\section{特征值的等价描述}

\[
\lambda\text{ 是 }T\text{ 的特征值}
\iff
\ker(T-\lambda I)\ne\{0\}.
\]

在有限维情形，还等价于：

\[
T-\lambda I\text{ 不可逆}.
\]

注意：

\[
T-\lambda I=0
\]

表示 $T=\lambda I$，即空间中每个向量都被同一个比例缩放。这比“$\lambda$ 是一个特征值”强得多。

\begin{jdefn}{特征空间}{r:4:2}
定义：

\[
E_\lambda
=
\ker(T-\lambda I).
\]

它称为 $\lambda$ 对应的特征空间。

其中的非零向量正是所有属于 $\lambda$ 的特征向量。
\end{jdefn}

\begin{jprop}{不同特征值对应的特征向量线性无关}{r:4:3}
若：

\[
\lambda_1,\ldots,\lambda_m
\]

两两不同，且：

\[
v_j\in E_{\lambda_j}\setminus\{0\},
\]

则：

\[
v_1,\ldots,v_m
\]

线性无关。
\end{jprop}

\begin{jproof}{证明思路}
假设它们存在最短的非平凡线性关系。

对该关系施加 $T-\lambda_m I$。最后一个向量会被消掉，而其余特征向量只会被乘上非零标量，因此得到一条更短的非平凡关系，产生矛盾。
\end{jproof}

\section{同一特征值也可能有多个线性无关特征向量}

若：

\[
\dim E_\lambda>1,
\]

则同一个特征值 $\lambda$ 对应多个线性无关特征向量。
